%% 第13回編集委員会（物理第1巻）
%% ph1-front.tex  前付の見開き2つ（物理シリーズ第1巻）
%%   p.4〜5  保存則と相似の地図（第8回の構成案「保存則とバランス式の表」と「相似の表」を図にしたもの）  \phConservationMapPages
%%   p.6〜7  この本の使い方（数学 bk-howto.tex の物理版。物理で足した囲みの見本をのせる）            \phHowToPages
%% 章のページ（\phchappage{力-2}）は ph1-pages.tex（gen_phys.py が作る）から読む。
\makeatletter
% 前付の見出し（数学の「この本の路線図」「この本の使い方」と同じ形）
\newcommand\ph@fronthead[2]{%
  \noindent{\sffamily\bfseries\fontsize{24}{30}\selectfont\color{bkDeep}#1}\hfill{\sffamily\fontsize{8.5}{11}\selectfont\color{bkGray}#2}\par\vspace{1mm}
  {\color{bkVol}\hrule height 1.4pt}\vspace{3.5mm}}
% 章の札（ページつき）。第1巻の章はページも出す。ほかの巻の章は札だけ
\newcommand\ph@cp[1]{\ph@routestop{#1}}
\newcommand\ph@cps[1]{\foreach \c [count=\k] in {#1} {\ifnum\k>1 \hspace{0.25em}\fi\ph@cp{\c}}}

%% ============================================================
%% p.4〜5  保存則と相似の地図
%% ============================================================
\newcommand\phConservationMapPages{%
\clearpage\setcounter{page}{4}%
\global\bk@decofalse\bkfolio{true}\bksetrunning{}{}{}%
{\def\R##1##2{##1}\pdfbookmark[0]{保存則と相似の地図}{ph-front-map}\typeout{BKPAGE front 保存則と相似の地図 \thepage}}%
\ph@fronthead{保存則と相似の地図}{全5巻の背骨}
{\fontsize{10.5}{19}\selectfont
物理のたくさんの法則の多くは、たった1つの形の式から生まれます。それが下の「バランス式」で、おこづかい帳のお金の出し入れと同じ形です。\par}
\vspace{1mm}
\begin{bkfigure}
\begin{tikzpicture}[x=1mm, y=1mm]
  % 見えない境界（検査体積）
  \draw[bkTheme, line width=1.1pt, dash pattern=on 3pt off 2pt, rounded corners=3mm, fill=bkTheme!5] (38,0) rectangle (102,27);
  \node[anchor=north west, text=bkDeep, font=\sffamily\bfseries\footnotesize] at (40,26) {見えない境界で囲んだ所};
  % たまる量（水位）
  \fill[bkBlue!22] (40,2) rectangle (100,13);
  \draw[bkBlue, line width=1pt] (40,13) -- (100,13);
  \node[text=bkBlue!70!black, font=\sffamily\bfseries\small] at (70,7.5) {たまる量（ふえる・へる）};
  % 生じる量
  \node[draw=phFlowR, line width=0.9pt, fill=white, star, star points=8, star point ratio=1.6, minimum size=6mm, inner sep=0pt] (g) at (55,18) {};
  \node[anchor=west, text=phFlowR!80!black, font=\sffamily\bfseries\small] at (60,18) {生じる量};
  % 入る・出る
  \draw[-{Stealth[length=3.6mm, width=3mm]}, line width=3.2pt, bkTheme] (6,15) -- (37,15);
  \node[anchor=south, text=bkDeep, font=\sffamily\bfseries\small] at (20,17) {入る量};
  \draw[-{Stealth[length=3.6mm, width=3mm]}, line width=3.2pt, bkGray] (103,15) -- (134,15);
  \node[anchor=south, text=bkInk, font=\sffamily\bfseries\small] at (119,17) {出る量};
\end{tikzpicture}
\end{bkfigure}
\vspace{1mm}
\begin{center}
\tikz\node[fill=bkTheme, text=white, rounded corners=2mm, inner xsep=5mm, inner ysep=2.6mm, font=\sffamily\bfseries\large]
  {入る量 \,$-$\, 出る量 \,$+$\, 生じる量 \,$=$\, たまる量};
\end{center}
\vspace{1mm}
{\fontsize{10.5}{19}\selectfont
「数える量」をかえると、力学・流体・熱・電気・原子炉の法則になります。生じる量が0の量は「保存する量」です。札の横の数は、この巻のページです。\par}
\vspace{2.5mm}
{\footnotesize\noindent
\begin{tblr}{colspec={Q[l,m,wd=19mm] Q[l,m,wd=25mm] Q[l,m,wd=17mm] X[l,m] Q[l,m,wd=36mm]}, width=\linewidth,
  row{1}={bg=bkVol!14, font=\sffamily\bfseries\footnotesize}, hlines={0.4pt,bkLine}, rowsep=0.5pt, colsep=2.5pt,
  column{1}={font=\sffamily\bfseries}}
数える量 & 入る・出るもの & 生じる量 & 法則・式 & 章 \\
質量 & 質量流量 & 0 & 連続の式 & \ph@cps{力-13} \\
運動量 & 力・運動量の流れ & 0 & 運動方程式、運動量保存、運動量の法則 & {\ph@cps{力-2,力-3}\\ \ph@cps{力-13}} \\
角運動量 & トルク & 0 & 回転の運動方程式、ポンプのオイラーの式 & \ph@cps{力-6,力-13} \\
エネルギー & 熱・仕事・エンタルピーの流れ・電磁場の流れ & 0 & 力学的エネルギー保存、第一法則、定常流動エネルギー式、ベルヌーイの式、ポインティングベクトル & {\ph@cps{力-4,力-13}\\ \ph@cps{熱-2,熱-4,電-14}} \\
電荷 & 電流 & 0 & 電荷保存、キルヒホッフの電流則 & \ph@cps{電-4} \\
エントロピー & 熱÷温度 & 0以上 & 第二法則 & \ph@cps{熱-6} \\
中性子 & 中性子の流れ & 生まれる数－吸収される数 & 中性子のバランス、増倍率 & \ph@cps{原-15} \\
\end{tblr}\par}
\clearpage
%% ---------------- p.5 相似 ----------------
\ph@fronthead{相似（にた形の式）の地図}{分野がちがっても、式は同じ形}
{\fontsize{10.5}{19}\selectfont
ばねにつけたおもりのゆれと回路の電流のゆれは、同じ形の式で書けます（一方を解けば、もう一方も解ける）。分野がちがっても式が同じ形になることを、\textbf{相似}とよびます。\par}
\vspace{1mm}
\begin{bkfigure}
\begin{tikzpicture}[x=1mm, y=1mm]
  \path[use as bounding box] (0,-17.5) rectangle (139,32);
  % 左：ばね・ダンパ・おもり
  \fill[bkLine!60] (0,4) rectangle (3,26);
  \draw[bkInk, line width=0.8pt] (3,4) -- (3,26);
  \foreach \y in {6,10,...,26} \draw[bkInk, line width=0.5pt] (0,\y) -- (3,\y-3);
  \draw[bkInk, line width=0.9pt, decorate, decoration={zigzag, amplitude=1.6mm, segment length=3mm}] (3,20) -- (31,20);
  \node[anchor=south, font=\sffamily\footnotesize] at (17,22.5) {ばね $k$};
  \draw[bkInk, line width=0.9pt] (3,9) -- (12,9) (22,5) -- (12,5) -- (12,13) -- (22,13) (17,9) -- (31,9);
  \draw[bkInk, line width=1.6pt] (17,6.3) -- (17,11.7);
  \node[anchor=north, font=\sffamily\footnotesize] at (17,4) {ダンパ $c$};
  \filldraw[fill=bkTheme!20, draw=bkTheme, line width=1pt, rounded corners=1mm] (31,4) rectangle (47,26);
  \node[font=\sffamily\bfseries\small] at (39,15) {$m$};
  \draw[-{Stealth[length=3mm]}, line width=1.4pt, phFlowR] (47,15) -- (60,15);
  \node[anchor=south, font=\sffamily\footnotesize] at (54,16.5) {力 $F$};
  \node[anchor=north, align=center, font=\small] at (31,-1) {$m\dfrac{\dd v}{\dd t}+c\,v+k\,x=F$\\[1mm]{\footnotesize $v=\dfrac{\dd x}{\dd t}$（速度）}};
  % まん中の等号
  \node[font=\sffamily\bfseries\Large, text=bkDeep] at (70,15) {＝};
  \node[font=\sffamily\bfseries\footnotesize, text=bkDeep] at (70,8) {同じ形};
  % 右：RLC の直列回路
  \begin{scope}[shift={(79,0)}]
    \draw[bkInk, line width=0.9pt] (6,4) -- (6,11) (6,19) -- (6,26) -- (16,26);
    \draw[bkInk, line width=0.9pt] (16,23.5) rectangle (28,28.5);
    \node[anchor=south, font=\sffamily\footnotesize] at (22,29) {抵抗 $R$};
    \draw[bkInk, line width=0.9pt] (28,26) -- (34,26);
    \draw[bkInk, line width=0.9pt] (34,26) \foreach \i in {0,...,3} { arc (180:0:1.6) } -- (52,26);
    \node[anchor=south, font=\sffamily\footnotesize] at (40.4,28.5) {コイル $L$};
    \draw[bkInk, line width=0.9pt] (52,26) -- (52,17.5) (49,17.5) -- (55,17.5) (49,15.5) -- (55,15.5) (52,15.5) -- (52,4) -- (6,4);
    \node[anchor=west, font=\sffamily\footnotesize] at (56,16.5) {$C$};
    \draw[bkInk, line width=0.9pt] (6,15) circle (4);
    \node[font=\footnotesize] at (6,15) {$e$};
    \draw[-{Stealth[length=2.6mm]}, line width=1.2pt, phFlowR] (20,6.5) -- (36,6.5);
    \node[anchor=north, font=\sffamily\footnotesize] at (28,5.2) {電流 $i$};
  \end{scope}
  \node[anchor=north, align=center, font=\small] at (111,-1) {$L\dfrac{\dd i}{\dd t}+R\,i+\dfrac{q}{C}=e$\\[1mm]{\footnotesize $i=\dfrac{\dd q}{\dd t}$（電流）}};
\end{tikzpicture}
\end{bkfigure}
\vspace{1mm}
{\footnotesize\noindent
\begin{tblr}{colspec={Q[l,m,wd=19mm] X[l,m] X[l,m] X[l,m] X[l,m]}, width=\linewidth,
  row{1}={bg=bkVol!14, font=\sffamily\bfseries\footnotesize}, hlines={0.4pt,bkLine}, rowsep=1pt, colsep=3pt,
  column{1}={font=\sffamily\bfseries}}
役わり & 力学（直線の動き） & 電気 & 熱 & 流体（管の流れ） \\
流れるもの & 速度 $v$ & 電流 $i$ & 熱流 & 体積流量 $Q$ \\
押すもの & 力 $F$ & 電圧 $e$ & 温度差 $\Delta T$ & 圧力差 $\Delta p$ \\
ためるもの & ばねのやわらかさ $1/k$ & 静電容量 $C$ & 熱容量 & タンク $A/(\rho g)$ \\
慣性 & 質量 $m$ & インダクタンス $L$ & （ない） & 管の中の流体 $\rho l/A$ \\
失うもの & ダンパ $c$ & 抵抗 $R$ & 熱抵抗 & 管路の抵抗 \\
\end{tblr}\par}
\vspace{1.5mm}
{\footnotesize\noindent\textbf{出会う章}\quad \ph@cp{力-2}\quad \ph@cp{力-6}\quad \ph@cp{力-14}\quad \ph@cp{波-1}（ゆれ）\quad \ph@cp{電-9}（回路）\quad 解き方は数学 6-12\par}
\begin{seikaku}
\begin{itemize}[leftmargin=1.2\zw, itemsep=1pt, topsep=1pt, label={\color{bkNavy}●}]
\item 力学の列は「力と電圧」、熱と流体の列は「差の量と電圧」が組です。だから質量・ばね・ダンパは直列の回路（上の図）、タンクは大地につないだコンデンサになります。
\item 熱流はすでに1秒あたりのエネルギーなので、温度差×熱流は仕事率になりません。
\item 乱流（力-14）の管路では、圧力差が流量の2乗に比例し、直線の関係になりません。
\item 力-15 の「相似則」（模型と本物の流れを同じ形にする条件）は、ここでの相似とは別の意味です。
\end{itemize}
\end{seikaku}
\clearpage\global\bk@decotrue}

%% ============================================================
%% p.6〜7  この本の使い方（物理の本）
%% ============================================================
\newcommand\ph@howstep[7]{% x, y, アイコン, 色, 名前, 説明, ノード名
  \node[draw=#4, line width=1pt, fill=white, rounded corners=2.5mm,
    minimum width=30.5mm, minimum height=33mm, anchor=north] (#7) at (#1,#2) {};
  \node[anchor=north] (#7i) at ([yshift=-2.2mm]#7.north) {\bkicon[1.7em]{#3}{#4}};
  \node[anchor=north, font=\sffamily\bfseries\small, text=#4, inner sep=0pt, align=center,
    minimum height=8.5mm, text width=29.5mm] (#7l) at ([yshift=-1mm]#7i.south) {#5};
  \node[anchor=north, font=\sffamily\footnotesize, text width=28mm, align=center, text=bkInk, inner sep=0pt]
    at ([yshift=-1.0mm]#7l.south) {#6};}
% 新しい囲みの見本の札  \ph@newbox{アイコン}{色}{名前}{説明}
\newcommand\ph@newbox[4]{%
  \begin{tikzpicture}[x=1mm, y=1mm]
    \draw[#2, line width=0.9pt, fill=white, rounded corners=2mm] (0,0) rectangle (67,14);
    \node[anchor=west, inner sep=0pt] (i) at (2.5,7) {\bkicon[1.5em]{#1}{#2}};
    \node[anchor=north west, inner sep=0pt, text=#2!75!black, font=\sffamily\bfseries\small] at (11.5,12) {#3};
    \node[anchor=north west, inner sep=0pt, text width=54mm, align=left, text=bkInk, font=\sffamily\fontsize{8.5}{12}\selectfont] at (11.5,6.6) {#4};
  \end{tikzpicture}}
\newcommand\phHowToPages{%
\clearpage\setcounter{page}{6}%
\global\bk@decofalse\bkfolio{true}\bksetrunning{}{}{}%
{\pdfbookmark[0]{この本の使い方}{ph-front-howto}\typeout{BKPAGE front この本の使い方 \thepage}}%
\ph@fronthead{この本の使い方}{}
{\fontsize{10.5}{19}\selectfont
この本は、身のまわりの「なぜ？」に物理で答える本です。できごとを、何かを無視した「モデル」にして式を立て、答えを数で出して確かめます。式は数学シリーズの道具で書きます。\par}

\bksubsection{1}{1つの章は、いつも同じ順番}
\begin{bkfigure}
\begin{tikzpicture}[x=1mm, y=1mm]
  \ph@howstep{0}{0}{hatena}{bkYellow!85!black}{はてな？}{身近なふしぎ\\から始めよう}{s1}
  \ph@howstep{35}{0}{yattemiyou}{bkGreen}{やってみよう}{実験して、\\はかってみよう}{s2}
  \ph@howstep{70}{0}{phmodel}{phModel}{モデル化の札}{何を無視するかを\\決めよう}{s3}
  \ph@howstep{105}{0}{phflow}{bkTheme}{一般の式から\\中高の公式へ}{条件で公式を導く}{s4}
  \ph@howstep{0}{-38}{reidai}{bkBlue}{例題}{7つの段で、\\ていねいに解く}{s5}
  \ph@howstep{35}{-38}{tsumazuki}{bkRed}{まちがえやすい\\ところ}{理由を知ろう}{s6}
  \ph@howstep{70}{-38}{phgear}{bkPurple}{工学ではこう使う}{本物の機械の\\数で確かめよう}{s7}
  \ph@howstep{105}{-38}{check}{bkTheme}{章末問題}{入試にチャレン\\ジもあるよ}{s8}
  \foreach \a/\b in {s1/s2,s2/s3,s3/s4,s5/s6,s6/s7,s7/s8}
    {\draw[bk arrow, bkGray, shorten <=1mm, shorten >=1mm] (\a.east) -- (\b.west);}
  \draw[bk arrow, bkGray, rounded corners=2mm, shorten <=1.5mm, shorten >=1.5mm] (s4.south) -- ++(0,-2.5) -| (s5.north);
\end{tikzpicture}
\end{bkfigure}
\vspace{1mm}
例題は「場面→何を無視するか→図と記号→使う式を選ぶ理由→式を立てる→解く→答えを確かめる」の7つの段で解きます。こまったら、例題にもどろう。

\bksubsection{2}{章のはじめの「なぜ？」と「道しるべ」}
\leavevmode 章の帯には、その章で答える「なぜ？」があります。帯のすぐ下には、こんな行があります。
\par\vspace{1.5mm}
\noindent\begin{minipage}{\linewidth}\bkmichishirube{4-11,5-5}{力-2,波-1}\end{minipage}
数字の札（4-11）は数学シリーズ、漢字の札（力-2）は物理シリーズの章です。わからない式は、数学の本にもどって確かめよう。

\clearpage
\bksubsection{3}{式が出てきたら、3つの見方で読む}
\begin{phyomi}[$v=gt$]
速さ $v$ は時間 $t$ に比例。1秒ごとに $g$（約\qty{9.8}{m/s}）ずつ速くなる。\par\noindent\phdim{m/s ＝ m/s$^2$ × s}\enspace{\small 左と右の単位がそろう}
\end{phyomi}
\vspace{-2mm}
\begin{phmijika}
ビルの3階（約\qty{10}{m}）から落ちると、約1.4秒後に約\qty{14}{m/s}（約\qty{50}{km/h}）。
\end{phmijika}
\vspace{-2mm}
\begin{phbaisuru}{時間 $t$}
速さも2倍。でも、落ちた距離は4倍（$t$ の2乗に比例する）。
\end{phbaisuru}

\bksubsection{4}{物理の本で足した囲み}
\noindent\ph@newbox{phmodel}{phModel}{モデル化の札}{無視したこと・仮定・くずれるとき}\hfill
\ph@newbox{phketa}{phKeta}{桁の感覚}{本物の値と、答えの桁の確かめ方}\par\vspace{2mm}
\noindent\ph@newbox{phstation}{phInk}{工学の駅}{この先の工学の科目と専門書のことば}\hfill
\ph@newbox{phflow}{bkTheme}{一般の式から中学・高校の公式へ}{条件をおいて、習った公式を導く}\par\vspace{2mm}
\noindent\ph@newbox{phmada}{phMada}{まだわかっていないこと}{今の科学でも答えが出ていないこと}\hfill
\ph@newbox{phsaki}{phSaki}{この先は大学（大学院）で}{この本の範囲をこえる理論の入口}\par

\bksubsection{5}{★と3つの段}
★は内容の高さ（★は必ず、☆の所まで）。★5の章は\textbf{目でわかる}→\textbf{手でわかる}→\textbf{式でわかる}の3段で、「ここまでで十分」の旗まで読めばゴールです。

\bksubsection{6}{この本の3つのやくそく}
\begin{enumerate}[label={\bkcircled[bkNavy]{\arabic*}}, leftmargin=2.4\zw, itemsep=2pt, topsep=2pt]
\item 式を立てる前に、図をかいて、何を無視するかを書こう。
\item 答えには単位をつけ、本物の値と桁が合っているかを確かめよう。
\item くり返しの計算は、表計算やプログラムにまかせよう（シミュレーション）。
\end{enumerate}
\clearpage\global\bk@decotrue}
\makeatother
